\scanpage{49}

II\hfill\mspage{- 10 -}\hfill n$^\circ$338

que $\displaystyle\sum_{i=1}^n x_ix'_i=1$. Les applications $A$-linéaires
$f:(a_i)\longrightarrow\displaystyle\sum_{i=1}^n a_ix_i$ de $A^n$
dans $J$ et $g:y\longrightarrow(x'_iy)$ de $J$ dans $A^n$ sont telles que
$f\circ g$ est l'application identique de $J$. Donc $g$ est un isomorphisme de $J$
sur un facteur direct du module libre $A^n$ (\uline{Alg.} VIII,\S1, n$^\circ$1 ), ce qui montre que $J$ est projectif.

Réciproquement supposons $J$ projectif. Il existe alors un isomorphisme $f$
de $J$ sur un facteur direct $M$ d'un $A$-module libre $A^{(I)}$ ; notons $g$ la surjection
de $A^{(I)}$ sur $J$ obtenue en composant la projection de $A^{(I)}$ sur $M$ avec
l'isomorphisme $f^{-1}$. Soit $(e_i)_{i\in I}$ la base canonique de $A^{(I)}$ ; posons
$x_i=g(e_i)$. L'application $\operatorname{pr}_i\circ f$ de $J$ dans $A$ est $A$-linéaire, et il existe
donc un élément $x'_i$ du corps des fractions $K$ de $A$ tel que $\operatorname{pr}_i(f(x))=x'_ix$
pour tout $x\in J$, (lemme 6). Pour $x\in J$ donné, les éléments $x'_ix$ ($i\in I$) sont
nuls sauf un nombre fini d'entre eux ; comme $J\ne(0)$, on peut prendre $x\ne0$,
ce qui montre que les $x'_i$ sont nuls sauf un nombre fini d'entre eux. D'autre
part, comme $g\circ f$ est l'application identique de $J$, on a
$x=\displaystyle\sum_{i\in I}xx'_i\,x_i$ pour
tout $x\in J$. Prenant $x\ne0$, on en déduit $1=\displaystyle\sum_{i\in I}x'_ix_i$.
Comme $g$ est surjectif, on a $J=\displaystyle\sum_{i\in I}Ax_i$.
D'autre part il résulte aussitôt de ce qu'on a vu
que l'idéal fractionnaire $J'=\displaystyle\sum_{i\in I}Ax'_i$ est l'inverse de $J$.
Ainsi $J$ est inversible, ce qui achève la démonstration.

Il résulte des prop. 5 et 3 qu'un $A$-module projectif de rang 1 est
nécessairement de type fini.

\margin{\unsure{?}}
Soit toujours $A$ un anneau
\figure{\begin{tikzpicture}[baseline=(w.base)]
\node[inner sep=0] (w) {intègre.};
\draw (.15,.12) .. controls (.15,.65) and (.7,.22) .. (.55,-.15)
 .. controls (.05,-.5) and (-.1,.1) .. (.15,.12);
\end{tikzpicture}}
L'ensemble $\underline I$ des idéaux fractionnaires
inversibles de $A$, muni de la multiplication, est un \uline{groupe commutatif} ;
l'ensemble $\underline P$ des idéaux fractionnaires principaux non nuls de $A$ est un sous-groupe
de $\underline I$. Le groupe quotient $\underline I/\underline P$ s'appelle le \uline{groupe des classes d'idéaux
de} $A$.

\margin{48}
